\chapter{Foreword}

\looseness=-1
Ajahn Karuṇadhammo has been part of the Abhayagiri \mbox{community} since its founding in 1996. Before joining the community, he worked as a nurse in Seattle, during which time he encountered the Theravadin tradition. Ajahn Karuṇadhammo received full bhikkhu ordination in 1998 under Luang Por Pasanno. Throughout his monastic life, he has served the community in many ways, including \mbox{as co-abbot.}

Over the years, Ajahn Karuṇadhammo has generously shared Dhamma reflections with the many people who have come to practice and listen to his talks. Offered with warmth and clarity, his teachings have encouraged many to pause, reflect, and gently turn the heart toward what is true.

The reflections gathered in this book are a selection based on Ajahn Karuṇadhammo's original talks, carefully edited for the written page while preserving their spirit and intention. These talks deeply resonated with us. As with other Dhamma teachings, they continue to reveal new layers of meaning through mindful reflection and practice.

\looseness=-1
For some time, it has been our aspiration to make these reflections available in written form. Most of Ajahn Karuṇadhammo's talks can be found online in audio form and continue to benefit listeners around the world. Reading, however, offers a different kind of engagement. It allows us to move at our own pace, to pause over a single sentence, to return to a passage that speaks to our present circumstances, or to revisit a reflection months or years later \mbox{with fresh understanding.}

Rather than presenting the Dhamma as a set of doctrines to be accepted, Ajahn Karuṇadhammo reminds us that ordinary experiences---moments of joy, disappointment, fear, loss, or confusion---can become the very ground for practice and opportunities for wisdom.

His reflections encourage us to look more carefully at the experience already unfolding in front of us. They remind us that liberation is cultivated not somewhere else, but in the midst of the very conditions of our daily lives. Rather than offering simple answers, these teachings often leave us with questions to carry into our own practice, inviting continued investigation and deeper understanding.

This book itself is also the fruit of Dhamma in action. During the past year, many conditions came together to make its publication possible. The project grew from Suwalee Sea Kochawat's aspiration to bring these reflections together in book form. I was fortunate to pick up the initial seeds from Sea and help coordinate the project through the transcribing and editing process. Jay Jacobs and Kristin Ohlson generously undertook the initial editing of the transcripts. Khantiko Bhikkhu offered further editing and proofreading, Hisayo Suzuki carefully attended to punctuation, copy editing, and proofreading, and Tyler Dennis provided typesetting support. We are also deeply grateful to everyone who offered financial support to make the printing \mbox{of this book possible.}

It has been deeply inspiring to witness how so many individuals, each contributing in their own way, helped bring this project to completion. This collaborative effort is a beautiful expression of dāna, generosity, and \mbox{spiritual friendship.}

We offer our heartfelt gratitude to Ajahn Karuṇadhammo for his years of patient teaching, his guidance, and his example, and to everyone whose generosity has made this book possible. May these reflections offer support and encouragement to all who open these pages, and may they help us return, again and again, to the direct knowing of the Dhamma, where true peace and freedom are found.

{\raggedleft
With gratitude,

  Anat Porat\\
  Suwalee Sea Kochawat\\
  October 2026
\par}
